\chapter{Conclusion}
\label{ch:conclusion}
The completed study connects a prepared BER aircraft-interval cohort to an existing object-centric airport simulation, exports interpretable logs and evaluates controlled timing, policy and resource scenarios. Attribution separates the supervisor's original process model from the student's empirical adapter, timing changes, provenance controls and analyses.

The 60-replication pilot completes 15,270 simulated turnaround executions based on 509 reused observed interval candidates. Literature timing assumptions shorten mean simulated block duration by approximately seven minutes when the inherited curfew helper is disabled. Much larger differences under inherited policy expose sensitivity to policy logic. Increasing fueling capacity removes most dispatch waiting but produces a much smaller overall turnaround change, consistent with overlap among parallel handling branches.

These findings provide partial answers to RQ1--RQ3. Available data inform demand and selected context, while unobserved services remain assumptions. Object-centric analysis exposes dependencies and extraction pitfalls, while stand-specific and inter-flight propagation comparisons remain open. Repeated capacity scenarios support conditional evaluation, not an optimal actual BER staffing plan.

The main practical requirement is to preserve the distinction between observation, assumption and generated behavior. Further work must correct the local-time movement policy, represent long stays and evaluate matching observable quantities on separate dates. This draft records the completed pilot and its limitations; it does not claim that the full empirical validation or thesis project is finished.
